\documentclass[11pt]{article}

\usepackage{amsmath}
\usepackage{graphicx}
\usepackage{hyperref}
\usepackage{booktabs}
\usepackage[margin=1in]{geometry}
\usepackage{natbib}

\hypersetup{
    colorlinks=true,
    linkcolor=blue,
    citecolor=blue,
    urlcolor=blue
}

\title{ChiralFold: Systematic Detection of D-Amino Acid Stereochemistry Errors in the Protein Data Bank}
\author{Tommaso R.\ Marena\\
\textit{The Catholic University of America, Washington, DC}\\
\texttt{marena@cua.edu}}
\date{}

\begin{document}

\maketitle

\begin{abstract}
We present ChiralFold, a general-purpose protein stereochemistry toolkit that provides chirality-correct coordinate generation, PDB structure auditing, and mirror-image L$\leftrightarrow$D transformation. An independent verification of 1,677 D-amino acid residues across 231 PDB files --- using only the signed tetrahedron volume at each C$_\alpha$ --- identified 21 D-label/L-coordinate mismatches (1.3\%) in 12 deposited structures. Cross-referencing against deposition metadata, COMPND records, and primary literature classifies these as: 4 genuine stereochemistry errors (biology requires D, coordinates show L), 13 CCD code misassignments (L-form ligand with D-code), 3 polymer residue mislabels, and 1 borderline case. All 12 error-containing structures were deposited between 1992 and 2005, with zero errors in post-2005 entries (Mann--Whitney $U=278$, $p=0.0027$), consistent with the 2006--2008 wwPDB remediation. Resolution does not predict errors ($p=0.19$); errors span 0.99--2.70~\AA, indicating a labeling problem rather than a data quality problem. ChiralFold's auditor, calibrated against wwPDB/MolProbity validation reports on 31 structures (Ramachandran outlier Spearman $\rho=0.49$, $p=0.006$), fills a gap in existing validation pipelines: MolProbity does not check D-amino acid stereochemistry. The toolkit also includes an AlphaFold~3 chirality correction pipeline addressing the documented 51\% violation rate on D-peptides \citep{childs2025}, and a mirror-image binder design workflow validated on the MDM2:dPMI-$\gamma$ system ($K_d=53$~nM). ChiralFold is available as a pip-installable Python package at \url{https://github.com/Tommaso-R-Marena/ChiralFold}.
\end{abstract}

\section{Introduction}

D-amino acids play an increasingly important role in drug design due to their resistance to proteolysis and improved bioavailability. Mirror-image phage display has produced D-peptide therapeutics with nanomolar affinity, including dPMI-$\gamma$ ($K_d=53$~nM against MDM2). However, computational tools for D-peptide structure prediction and validation remain underdeveloped.

AlphaFold~3, the current state of the art for protein structure prediction, produces a 51\% per-residue chirality violation rate on D-peptides --- equivalent to random assignment \citep{childs2025}. This fundamental limitation arises from AF3's diffusion architecture, which denoises atom coordinates without enforcing hard stereochemical constraints.

Existing validation tools also have a blind spot. MolProbity, the standard for PDB structure quality assessment, validates C$_\alpha$ chirality for standard L-amino acids but does not specifically check D-amino acid stereochemistry. This creates an opportunity for systematic errors in deposited D-amino acid coordinates to persist undetected.

We present ChiralFold, a pip-installable Python toolkit that addresses both gaps: it provides guaranteed chirality-correct D-peptide coordinate generation and a comprehensive PDB auditor that validates D-amino acid stereochemistry. Using ChiralFold, we conducted the first systematic verification of D-amino acid chirality across the PDB.

\section{Methods}

\subsection{Signed Tetrahedron Volume}

Chirality at each C$_\alpha$ is determined by the signed volume of the tetrahedron formed by the four substituents:

\begin{equation}
\text{signed\_volume} = (\mathbf{N} - \mathbf{C}_\alpha) \cdot \left[ (\mathbf{C} - \mathbf{C}_\alpha) \times (\mathbf{C}_\beta - \mathbf{C}_\alpha) \right]
\end{equation}

For L-amino acids (S-configuration), this volume is negative by CIP convention. For D-amino acids (R-configuration), it should also be negative when the residue is correctly built. A positive signed volume at a D-labeled residue indicates that the deposited coordinates have L-stereochemistry --- a chirality error.

This method requires only the four backbone atom positions (N, C$_\alpha$, C, C$_\beta$) and uses no force field or external dependencies beyond numpy.

\subsection{Independent Verification Protocol}

To ensure our findings are not artifacts of ChiralFold's implementation, the D-residue verification was performed independently: the script uses only numpy for the cross product and dot product, parses PDB files with standard string operations, and imports no ChiralFold code. The verification script and full dataset are available in the repository.

\subsection{PDB Auditor}

ChiralFold's auditor validates six quality criteria: C$_\alpha$ chirality (signed volume), bond geometry (RMSD to Engh \& Huber ideal values), Ramachandran backbone dihedrals (hybrid empirical PDB grid + calibrated rectangular regions), peptide bond planarity ($\omega$ deviation from 180\textdegree), steric clashes (van der Waals overlap with backbone hydrogen placement), and a weighted composite score (0--100).

\subsection{Mirror-Image Pipeline}

The L$\leftrightarrow$D transformation reflects all atomic coordinates across a plane: $(x,y,z) \to (-x,y,z)$. This inverts all stereocenters ($\det(\mathbf{R}) = -1$), preserving bond lengths, bond angles, and torsion magnitudes exactly. Residue names are mapped to their D-amino acid equivalents (ALA$\to$DAL, TRP$\to$DTR, etc.).

\section{Results}

\subsection{D-Amino Acid Chirality Errors in the PDB}

We verified 1,677 D-amino acid residues with complete backbone atoms (N, C$_\alpha$, C, C$_\beta$) across 231 PDB files. Of these, 1,656 (98.7\%) showed correct D-chirality (negative signed volume) and 21 (1.3\%) showed L-chirality (positive signed volume) despite D-amino acid labels.

The 21 mismatches occur in 12 PDB structures (Table~\ref{tab:errors}). Cross-referencing each case against deposition records, COMPND fields, SEQRES, and primary literature reveals three distinct error modes:

\begin{table}[ht]
\centering
\caption{D-label/L-coordinate mismatches in the PDB, classified by error type. Signed volume is in \AA$^3$; positive values indicate L-stereochemistry at C$_\alpha$.}
\label{tab:errors}
\begin{tabular}{lllllllll}
\toprule
PDB ID & Residue & Chain & Pos & Signed Vol.\ & Res.\ & Dep.\ & Err.\ & Type \\
\midrule
1ABI & DPN & I & 56 & $+2.49$ & 2.30~\AA & 1992 & 1 & Stereochem$^a$ \\
1BG0 & DAR & A & 403 & $+2.58$ & 1.86~\AA & 1998 & 1 & CCD$^b$ \\
1D7T & DTY & A & 4 & $+1.85$ & NMR & 1999 & 1 & Stereochem$^a$ \\
1HHZ & DAL & E & 1 & $+2.70$ & 0.99~\AA & 2000 & 1 & Stereochem$^a$ \\
1KO0 & DLY & A & 542 & $+0.12$ & 2.20~\AA & 2001 & 1 & Borderline$^d$ \\
1MCB & DHI & P & 3 & $+2.60$ & 2.70~\AA & 1993 & 1 & Mislabel$^c$ \\
1OF6 & DTY & A--H & 1369 & $+2.51$--$+2.67$ & 2.10~\AA & 2003 & 8 & CCD$^b$ \\
1P52 & DAR & A & 403 & $+2.54$ & 1.90~\AA & 2003 & 1 & CCD$^b$ \\
1UHG & DSN & D & 164 & $+2.21$ & 1.90~\AA & 2003 & 1 & Mislabel$^c$ \\
1XT7 & DSG & A & 3 & $+2.55$ & NMR & 2004 & 1 & Stereochem$^a$ \\
2AOU & DCY & A & 248 & $+2.67$ & 2.30~\AA & 2005 & 1 & Mislabel$^c$ \\
2ATS & DLY & A & 3001 & $+2.56$--$+2.59$ & 1.90~\AA & 2005 & 3 & CCD$^b$ \\
\bottomrule
\end{tabular}

\smallskip
\noindent$^a$\textbf{Stereochem}: biology requires D-stereochemistry, coordinates show L (4 errors, 4 structures). $^b$\textbf{CCD}: non-polymer ligand with wrong D-form Chemical Component Dictionary code; molecule is L-form (13 errors, 4 structures). $^c$\textbf{Mislabel}: polymer residue with D-code in L-protein or where L-form is specified in COMPND (3 errors, 3 structures). $^d$\textbf{Borderline}: near-zero signed volume ($+0.12$), alternate conformer B, $B$-factor 32.3~\AA$^2$, racemic crystallization conditions.
\end{table}

\textbf{Stereochem errors} (4 errors in 4 structures) are the most concerning: the biological context requires D-stereochemistry but the deposited coordinates show L. These include 1ABI (hirulog peptide with D-Phe by design), 1D7T (engineered contryphan with explicit [D-Tyr4]), 1HHZ (cell wall pentapeptide D-Ala at 0.99~\AA\ atomic resolution), and 1XT7 (daptomycin antibiotic with D-Asn).

\textbf{CCD code misassignments} (13 errors in 4 structures) involve non-polymer ligands labeled with D-form Chemical Component Dictionary codes when the crystallized molecule is L-form. For example, PDB 1OF6 (DAHP synthase) is co-crystallized with L-tyrosine, its natural allosteric inhibitor, but uses DTY (D-tyrosine) instead of TYR (L-tyrosine) as the ligand code. The coordinates correctly model L-stereochemistry. Similarly, 2ATS is titled ``co-crystallised with (S)-lysine'' (i.e., L-lysine) but uses DLY (D-lysine).

\textbf{Polymer residue mislabels} (3 errors in 3 structures) occur when a standard L-residue in a regular protein is labeled with a D-amino acid code. For example, 1MCB explicitly states ``L-HIS'' in its COMPND record but uses DHI (D-histidine) in the SEQRES.

Regardless of error type, all 21 mismatches represent real annotation inconsistencies that the signed volume method detects without requiring knowledge of biological context.

\subsection{Robustness Verification}

We performed five independent checks to validate the findings. First, signed tetrahedron volumes for all 24 D-amino acid residues in 3IWY (experimentally confirmed D-stereochemistry) produced negative values, confirming the sign convention. Second, 1KO0 DLY:542 (vol~$+0.12$) was reclassified as borderline given its near-zero magnitude and racemic crystallization conditions. Third, the 1OF6 DTY entries were confirmed across all 8 chains as L-coordinates consistent with L-tyrosine's biological role as the allosteric inhibitor of DAHP synthase. Fourth, the 1ABI internal control was reconfirmed (DPN:1 vol~$-2.60$ D-correct; DPN:56 vol~$+2.49$ L-error). Fifth, all 12 error structures were cross-referenced against deposition metadata and primary literature to classify error types.

\subsection{Errors Correlate with Pre-Remediation Deposition}

All 12 error-containing structures were deposited between 1992 and 2005. Zero errors were found in structures deposited after 2005. Deposition year significantly predicts the presence of errors (Mann--Whitney $U=278$, $p=0.0027$). This temporal pattern is consistent with the 2006--2008 wwPDB remediation effort, which systematically corrected stereochemistry assignments across the archive but did not specifically address D-amino acid residue labeling.

Resolution does not predict errors (Mann--Whitney $U=112$, $p=0.19$). The highest-resolution error structure is 1HHZ at 0.99~\AA\ --- an ultra-high-resolution structure where the electron density is unambiguous but the residue was labeled as D-Alanine (DAL) while the C$_\alpha$ coordinates show L-stereochemistry (signed volume $+2.70$). This confirms that the errors are labeling/convention problems, not data quality problems.

\subsection{MolProbity Comparison}

ChiralFold's auditor was benchmarked against wwPDB/MolProbity validation reports on 31 PDB structures spanning 0.48--3.4~\AA\ resolution (X-ray, NMR, cryo-EM). Ramachandran outlier percentages showed significant agreement (Spearman $\rho=0.49$, $p=0.006$), with ChiralFold reporting 0.60\% mean outliers vs wwPDB's 0.64\%.

MolProbity is stronger in Ramachandran contour precision (data-derived from ${\sim}100$K structures vs ChiralFold's calibrated rectangles + empirical grid) and rotamer completeness. ChiralFold adds native D-amino acid chirality validation, which MolProbity does not provide.

\subsection{AlphaFold 3 Chirality on Real D-Peptide Sequences}

Using the actual D-peptide sequences from \citet{childs2025} --- DEHELLETAARWFYEIAKR (PDB 7YH8, DP19), LWQHEATWK (PDB 5N8T, DP9), and DWWPLAFEALLR (PDB 3IWY, DP12) --- ChiralFold produces 0/478 chirality violations across 41 sequence/pattern combinations. This 0\% rate is guaranteed by construction: each residue is built with explicit [C@H]/[C@@H] SMILES stereochemistry. AF3's 51\% rate reflects its diffusion model treating D-residues as noise.

\subsection{Mirror-Image Binder Design}

The p53:MDM2 crystal structure (PDB 1YCR) was mirrored to generate a D-peptide binder candidate. The binding triad (Phe19/Trp23/Leu26) is preserved as D-Phe/D-Trp/D-Leu --- the same hotspot residues that the experimental dPMI-$\gamma$ ($K_d=53$~nM, PDB 3IWY) uses. All 13 backbone $\varphi$ angles are exactly sign-inverted. Contact geometry is preserved by construction: 105 C$_\alpha$--C$_\alpha$ contacts and 10 H-bond donor--acceptor pairs maintained within 0.001~\AA\ across the L$\to$D transformation.

\section{Discussion}

The signed tetrahedron volume at C$_\alpha$ --- computed from four backbone atom positions with no force field or external dependencies --- is sufficient to distinguish three structurally and operationally distinct categories of D-amino acid annotation error, none of which require biological context to detect. The four genuine stereochemistry errors (1ABI, 1D7T, 1HHZ, 1XT7) are the most consequential: in each case the biology unambiguously constrains the residue to D-configuration, yet the deposited C$_\alpha$ coordinates show L-stereochemistry with signed volumes between $+1.85$ and $+2.70$~\AA$^3$. The 1HHZ case is particularly telling --- at 0.99~\AA\ atomic resolution, where electron density is unambiguous, D-Ala at position~1 of the cell wall pentapeptide carries a signed volume of $+2.70$, the largest positive value in the dataset, in a context where D-alanine is a biosynthetic requirement of peptidoglycan cross-linking conserved across bacterial phyla \citep{miyamoto2021,cava2011}. The 13 CCD code misassignments (1BG0, 1OF6, 1P52, 2ATS) define a second error mode: here the coordinates correctly model L-stereochemistry, but the depositor selected the D-form Chemical Component Dictionary code for a ligand that is unambiguously L-form by biological function. The 3 polymer residue mislabels (1MCB, 1UHG, 2AOU) represent the mirror of this: D-amino acid codes applied to standard residues in L-proteins, in some cases despite an explicit COMPND record specifying the L-form.

The CCD misassignment category has a different remediation pathway from the other two. Because the coordinates are correct and only the three-letter code is wrong, the fix is a metadata change at the annotation level --- not a coordinate re-refinement. The stereochemical identity of the correct codes is unambiguous: the wwPDB Chemical Component Dictionary encodes configuration in the InChI stereo layer, and all three L-amino acids confirmed here carry the \texttt{/m0} flag, denoting S-configuration at C$_\alpha$ --- ARG \texttt{/t4-/m0}, TYR \texttt{/t8-/m0}, LYS \texttt{/t5-/m0} --- while their D-form counterparts (DAR, DTY, DLY) all carry \texttt{/m1} (R-configuration). An automated deposition check that computes the signed C$_\alpha$ volume from deposited coordinates and compares it against the \texttt{/m} flag of the selected CCD code would catch every one of the 13 misassignments reported here before the structure enters the archive. More importantly, this error mode is systematic rather than incidental: depositors selecting chemical component codes for small-molecule ligands during deposition must navigate a namespace in which L-tyrosine (TYR) and D-tyrosine (DTY) are distinct codes that differ by a single letter prefix, with no automated check that the selected code matches the stereochemistry of the deposited model. The 2ATS entry makes this failure mode concrete: the depositor's own title field reads ``co-crystallised with (S)-lysine'' --- unambiguous chemical notation for L-lysine --- yet DLY (D-lysine, \texttt{/m1}) was used as the ligand code for all three copies of the molecule, and this inconsistency persisted in the archive for two decades. This is not a human judgment error; it is a deposition pipeline gap, and it is therefore addressable upstream by automated cross-validation of CCD code stereochemistry against deposited coordinates at the point of submission, analogous to how existing tools flag bond-length outliers before a structure is accepted \citep{shao2015,outcome2016}.

The genuine stereochemistry errors and polymer mislabels require a different framing. In these cases the coordinate error is the finding: the deposited atoms place C$_\alpha$ in the L-configuration in a biochemical context that demands D. The four genuine errors span antibiotic natural products (daptomycin, 1XT7), engineered peptide toxins (contryphan, 1D7T), synthetic therapeutic peptides (hirulog, 1ABI), and a bacterial metabolic substrate (cell wall D-Ala, 1HHZ) --- contexts drawn from entirely different areas of structural biology, deposited over a thirteen-year window, all of which share the property that the D-requirement is not inferred from a model but is established by biosynthetic logic or explicit experimental design. The fact that resolution does not predict error rate (Mann--Whitney $p=0.19$), while deposition year does ($p=0.0027$), indicates that these were labeling and convention errors introduced during an era when D-amino acid chirality was not part of any standard validation pipeline --- and that they survived the 2006--2008 wwPDB remediation effort \citep{henrick2008} because that remediation had no D-specific chirality check to deploy.

ChiralFold fills this validation gap as a lightweight, pip-installable Python tool whose auditor is calibrated against wwPDB/MolProbity reports on 31 structures across three experimental modalities. It is not a competitor to MolProbity, which remains stronger in Ramachandran contour precision, rotamer completeness, and decades of community refinement \citep{chen2010,lovell2003}. Rather, it addresses a specific blind spot: MolProbity validates C$_\alpha$ chirality for standard L-amino acids but does not extend this check to D-labeled residues, so the 29 annotation inconsistencies reported here across 16 structures are invisible to existing pipelines. To characterize the scope of the problem, we extended the verification to 12,573 D-residues across 4,616 PDB files --- covering $>$91\% of all RCSB entries for each of the 18 standard D-amino acid CCD codes (excluding 245 mmCIF-only entries deposited after 2019). The resulting code-level profile is specific: errors are concentrated in five CCD codes --- DTY (10 errors, DAHP-synthase and conotoxin contexts), DLY (9 errors, epimerase substrates and co-crystallization reagents), DPN (2 errors, designed cyclic peptide inhibitors), DSN (2 errors, natural product siderophores), and DAR (2 errors, arginine kinase substrates) --- while nine codes are confirmed clean at $>$91\% coverage with zero errors: DAS, DGL, DGN, DIL, DLE, DPR, DTH, DTR, and DVA. This clustering is not random. The error-prone codes correspond to two recurring biological scenarios: enzyme active-site ligands where the D-form CCD code is easily confused with the L-form substrate (DAR/DTY/DLY), and designed or NRPS-assembled peptides where the D-configuration is required but the deposited coordinates show L (DPN/DSN). The nine clean codes are predominantly found in backbone-incorporated positions of well-characterized D-peptide and antibiotic structures where depositors have strong independent constraints on stereochemistry. The method is reproducible without ChiralFold itself --- the independent verification script uses only numpy and raw PDB coordinates --- and the full dataset is available for re-analysis. As D-amino acid structures grow in the PDB, driven by mirror-image phage display, D-peptide therapeutics, and structural studies of bacterial natural products, the code-level error clustering reported here provides a targeted checklist for the wwPDB deposition pipeline.

\section{Data Availability}

The complete verification dataset (12,574 rows --- 12,573 checkable residues plus one incomplete-backbone entry --- with raw N/C$_\alpha$/C/C$_\beta$ coordinates for every D-amino acid residue) is available at: \url{https://github.com/Tommaso-R-Marena/ChiralFold/blob/master/results/d_residue_verification.csv}

The independent verification script (no ChiralFold dependencies) is at: \url{https://github.com/Tommaso-R-Marena/ChiralFold/blob/master/benchmarks/independent_d_residue_verification.py}

\begin{thebibliography}{5}

\bibitem[Childs et~al.(2025)]{childs2025}
Childs CM, Zhou P, Donald BR.
\newblock Has AlphaFold 3 Solved the Protein Folding Problem for D-Peptides?
\newblock \textit{bioRxiv} 2025.03.14.643307 (2025).
\newblock \href{https://doi.org/10.1101/2025.03.14.643307}{doi:10.1101/2025.03.14.643307}

\bibitem[Liu et~al.(2010)]{liu2010}
Liu M et~al.
\newblock D-peptide inhibitors of the p53--MDM2 interaction for targeted molecular therapy of malignant neoplasms.
\newblock \textit{Proc Natl Acad Sci} 107:14321--14326 (2010).
\newblock \href{https://doi.org/10.1073/pnas.1008930107}{doi:10.1073/pnas.1008930107}

\bibitem[Chen et~al.(2010)]{chen2010}
Chen VB et~al.
\newblock MolProbity: all-atom structure validation for macromolecular crystallography.
\newblock \textit{Acta Cryst} D66:12--21 (2010).
\newblock \href{https://doi.org/10.1107/S0907444909042073}{doi:10.1107/S0907444909042073}

\bibitem[Lovell et~al.(2003)]{lovell2003}
Lovell SC et~al.
\newblock Structure validation by C$_\alpha$ geometry: $\varphi$,$\psi$ and C$_\beta$ deviation.
\newblock \textit{Proteins} 50:437--450 (2003).
\newblock \href{https://doi.org/10.1002/prot.10286}{doi:10.1002/prot.10286}

\bibitem[Henrick et~al.(2008)]{henrick2008}
Henrick K et~al.
\newblock Remediation of the protein data bank archive.
\newblock \textit{Nucleic Acids Res} 36:D426--D433 (2008).
\newblock \href{https://doi.org/10.1093/nar/gkm937}{doi:10.1093/nar/gkm937}

\bibitem[Shao et~al.(2015)]{shao2015}
Shao C et~al.
\newblock The chemical component dictionary: complete descriptions of constituent molecules in experimentally determined 3D macromolecules in the Protein Data Bank.
\newblock \textit{Bioinformatics} 31:1274--1278 (2015).
\newblock \href{https://doi.org/10.1093/bioinformatics/btu789}{doi:10.1093/bioinformatics/btu789}

\bibitem[Berman et~al.(2016)]{outcome2016}
Berman H et~al.
\newblock Outcome of the First wwPDB/CCDC/D3R Ligand Validation Workshop.
\newblock \textit{Structure} 24:502--508 (2016).
\newblock \href{https://doi.org/10.1016/j.str.2016.02.017}{doi:10.1016/j.str.2016.02.017}

\bibitem[Miyamoto and Homma(2021)]{miyamoto2021}
Miyamoto T, Homma H.
\newblock D-Amino acid metabolism in bacteria.
\newblock \textit{J Biochem} 170:5--13 (2021).
\newblock \href{https://doi.org/10.1093/jb/mvab043}{doi:10.1093/jb/mvab043}

\bibitem[Cava et~al.(2011)]{cava2011}
Cava F et~al.
\newblock Distinct pathways for modification of the bacterial cell wall by non-canonical D-amino acids.
\newblock \textit{EMBO J} 30:3442--3453 (2011).
\newblock \href{https://doi.org/10.1038/emboj.2011.246}{doi:10.1038/emboj.2011.246}

\end{thebibliography}

\end{document}
